\section{The Haim--Kislev Quadratic Program}
\label{sec:quadratic-program-algorithm-hk2019}

This section records the finite reduction used later in the thesis.
Haim--Kislev's formula for polytopes reduces the EHZ capacity to a finite list
of quadratic programs over ordered facet words
\cite{HK2017}. For general polytopes this list is combinatorially large. For
Lagrangian products of planar polygons, the billiard interpretation gives a
much smaller finite list; this finite-enumeration result is the part needed in
Section~\ref{sec:rotated-regular-polygons}.
After this first attribution, we call the resulting finite optimization problem
the quadratic program, or QP when space is tight.

\subsection{The finite quadratic program}
\label{subsec:quadratic-program-algorithm-hk2019-general-polytopes}

Let
\[
  K=\{x\in\mathbb{R}^4:\langle a_i,x\rangle\leq 1,\;
    i=1,\ldots,F\}
\]
be a full-dimensional polytope with \(0\in\operatorname{int}(K)\), written as
in Section~\ref{sec:generalized-reeb-orbits-polytopes}. Thus
\(a_i=n_i/h_i\) if \(n_i\) is the outward unit normal to the \(i\)-th facet and
\(h_i=h_K(n_i)>0\) is its height. Haim--Kislev state their formula in the
\((n_i,h_i)\)-variables. We use the equivalent dual-vertex variables
\[
  \beta_i^{\mathrm{here}}=\beta_i^{\mathrm{HK}}h_i.
\]
The constraints then become
\[
  \beta_i\geq0,\qquad
  \sum_i\beta_i=1,\qquad
  \sum_i\beta_i a_i=0.
\]

For an ordered facet word
\(\sigma=(\sigma_1,\ldots,\sigma_m)\) with distinct entries, define
\[
  M_\sigma
  =
  \left\{
    \beta\in\mathbb{R}_{\geq0}^m:
    \sum_{k=1}^m\beta_k=1,\;
    \sum_{k=1}^m\beta_k a_{\sigma_k}=0
  \right\}.
\]
The quadratic objective for this active word is
\begin{equation}\label{eq:quadratic-program-algorithm-hk2019-q}
  Q_\sigma(\beta)
  =
  \sum_{1\leq j<k\leq m}
    \beta_j\beta_k\,\omega_0(a_{\sigma_j},a_{\sigma_k}).
\end{equation}
Here \(\sigma\) is written in the Reeb traversal order used in
Section~\ref{sec:generalized-reeb-orbits-polytopes}. Haim--Kislev's displayed
formula uses the reversed ordered word after translating their symplectic
convention to \(\omega_0\). Since word reversal is a bijection on the ordered
candidate family, the maximum is unchanged; the convention matters only when a
single fixed word is interpreted.

Let
\[
  Q_{\max}
  =
  \max_\sigma\ \max_{\beta\in M_\sigma} Q_\sigma(\beta),
\]
where \(\sigma\) ranges over ordered words with distinct facets. Equivalently,
one may range over all permutations of all facets and allow zero weights; the
zero-weight entries are then deleted to recover the active word. This is only
an enumeration convention: HK's theorem is the full-permutation statement, and
the active-word form is obtained by deleting unused zero-weight facets and
reindexing the remaining entries. This formulation intentionally permits
duplicates of the same geometric orbit, for example from changing the starting
segment.

\begin{theorem}[Haim--Kislev quadratic program]
\label{thm:quadratic-program-algorithm-hk2019-hk-formula}
For a full-dimensional convex polytope \(K\subset\mathbb{R}^4\) with
\(0\in\operatorname{int}(K)\), written in the facet normalization above,
\[
  c_{\mathrm{EHZ}}(K)=\frac{1}{2Q_{\max}}.
\]
\end{theorem}

\begin{proof}
This is Haim--Kislev's finite capacity formula
\cite[Theorem~1.1]{HK2017}, rewritten from unit normals and heights to the
dual vertices \(a_i=n_i/h_i\). The variable change above sends
\[
  \sum_i\beta_i^{\mathrm{HK}}h_i=1,\qquad
  \sum_i\beta_i^{\mathrm{HK}}n_i=0
\]
to the constraints defining \(M_\sigma\). Their displayed objective uses the
opposite word orientation from \eqref{eq:quadratic-program-algorithm-hk2019-q};
replacing each ordered word by its reverse sends it to \(Q_\sigma\). Since
reversal is a bijection on the words in the maximum, their formula reads
\[
  c_{\mathrm{EHZ}}(K)
  =
  \frac{1}{2}
  \left[\max_\sigma\max_{\beta\in M_\sigma}Q_\sigma(\beta)\right]^{-1}.
\]
\end{proof}

For a fixed active word, the KKT system is a stationarity condition on the
relative interior of the corresponding support face.  A solution with
positive coordinates is therefore an exact feasible stationary witness; it is
not, by itself, a certificate that this fixed word maximizes~$Q_\sigma$, nor
that it represents a Reeb orbit.  This distinction does not prevent the finite
Haim--Kislev algorithm from computing the scalar capacity.  Every positive
witness is feasible and hence has value at most~$Q_{\max}$, while a global
maximizer becomes a positive KKT solution after its zero coordinates are
deleted.  Thus exact KKT solutions suffice when the complete candidate family
is enumerated.  A claim about a particular maximizing word, a physical orbit,
or a complete minimizer set cannot be inferred from an arbitrary KKT witness.
There is, however, no further gap for a \emph{global} QP maximizer: the
construction in the proof of Haim--Kislev's formula turns it into a global
minimizer of the dual problem, and the Clarke correspondence of
Theorem~\ref{thm:preliminaries-clarke-dual-action} supplies a translation and
positive rescaling that realize it as a minimum-action generalized Reeb orbit.
After deleting zero-weight entries, its pure facet directions occur in the
active order~$\sigma$, so the realized orbit is simple.  Feasibility or KKT
stationarity without global optimality does not provide this conclusion.

The normalization is compatible with our Reeb-orbit convention. Suppose a
simple orbit is traversed in the active order
\(\sigma=(\sigma_1,\ldots,\sigma_m)\), with dwell times \(\tau_k>0\) and period
\(T=\sum_k\tau_k\). Set
\[
  \beta_k=\frac{\tau_k}{T}.
\]
Then \(\beta\in M_\sigma\), because the closure equation
\(\sum_k\tau_kR_{\sigma_k}=0\) and \(R_i=2J_0a_i\) imply
\(\sum_k\beta_k a_{\sigma_k}=0\). Lemma~\ref{lem:preliminaries-shoelace}
gives
\[
  2T
  =
  4T^2
  \sum_{1\leq r<s\leq m}
    \frac{\tau_r}{T}\frac{\tau_s}{T}\,
    \omega_0(a_{\sigma_r},a_{\sigma_s})
  =
  4T^2Q_\sigma(\beta),
  \qquad\text{so}\qquad
  Q_\sigma(\beta)=\frac{1}{2T}.
\]
Thus maximizing \(Q_\sigma\) is the same normalization as minimizing the
free-period action \(T=A(\gamma)\). The Haim--Kislev theorem is the part that
rules out a larger value coming from the finite optimization and proves the
global equality with the capacity.

For a fixed active word, it is often convenient to write the objective as a
quadratic program
\[
  \text{maximize }\frac{1}{2}\beta^{\mathsf T}H_\sigma\beta
  \quad\text{subject to}\quad
  \beta\in M_\sigma,
\]
where \(H_\sigma\) is the symmetric matrix determined by
\(\frac{1}{2}\beta^{\mathsf T}H_\sigma\beta=Q_\sigma(\beta)\). This section
only fixes the standard finite problem. The next subsection explains how later
finite certificates use this fixed problem.

\subsection{From the formula to finite checks}
\label{subsec:quadratic-program-algorithm-hk2019-computation-layers}

The Haim--Kislev theorem supplies the finite optimization problem, but it does
not prescribe a single proof layout for later finite checks. In this thesis,
later arguments separate three choices: which words are still under
consideration, what is checked for one fixed active word, and which exact
assertions are needed to turn the finite list into a theorem.

The reference word family is the unpruned Haim--Kislev family: all ordered
cyclic words with distinct facets. This is equivalent to the all-facet
permutation formulation in Theorem~\ref{thm:quadratic-program-algorithm-hk2019-hk-formula},
because a zero component of \(\beta\) contributes nothing to the constraints or
to \(Q_\sigma\), and deleting it gives the corresponding shorter active word.
Conversely, an active word can be padded by unused facets with zero weight.

For a fixed active word, a capacity-relevant simple orbit has \(\beta_k>0\) for
each visited facet. If the maximum for a longer word is attained only with some
\(\beta_k=0\), then the same value is represented by the shorter active word
obtained by deleting that facet. Thus later fixed-word checks use the interior
KKT equations when they claim a smooth branch for a chosen active word.
Boundary maximizers belong to shorter active words, not to new zero-dwell
orbits.

The ordinary general-polytope search may prune the reference word family by
directed transition tests. A simple orbit that moves from facet \(F_i\) to
facet \(F_j\) must have a breakpoint in \(F_i\cap F_j\), and the direction of
the move imposes the corresponding sign condition on
\(\omega_0(a_i,a_j)\). Stronger pruning may add other necessary transition
checks, but every discarded pair must fail a necessary condition for that
directed transition. Passing such tests is only a necessary condition for a
physical transition; it is not itself a proof that the whole word has a valid
base point. The fixed active-word solve and, when needed, the base-point
recovery check from Lemma~\ref{lem:generalized-reeb-base-point-recovery}
provide that later validation for a non-global candidate.  For a certified
global QP maximizer, orbit existence and primal placement instead follow
globally from the dual-minimizer/Clarke argument above; base-point recovery is
then a finite reconstruction method, not an additional existence hypothesis.

For Lagrangian products we use a different enumeration route, described in
Theorem~\ref{thm:lagrangian-product-finite-enumeration}. It feeds the same
active-word quadratic programs but visits only the alternating block family that
contains a capacity minimizer for such products.

The active-word solve is a quadratic program, equivalently a KKT calculation.
For theorem-bearing finite certificates, the load-bearing claims are exact
algebraic assertions: feasibility, positivity, branch formulas, sign decisions,
rank statements, or inequalities. Numerical searches may suggest the finite
list or the expected branch, but they are not used as proof unless replaced by
such exact assertions. The numerical issues are discussed in
Section~\ref{sec:numerics}.  The exact certificates and their trust boundaries are
described with the HKO argument in
Section~\ref{subsec:hko-local-maximum-computation-with-sagemath} and with the
rotated-pentagon argument in Section~\ref{sec:rotated-regular-polygons}.

\subsection{Finite enumeration for Lagrangian products}
\label{subsec:quadratic-program-algorithm-hk2019-lagrangian-products}

Let \(K_q,K_p\subset\mathbb{R}^2\) be nondegenerate convex polygons containing
the origin in their interiors, and let
\[
  K=K_q\times_L K_p\subset\mathbb{R}^2_q\oplus\mathbb{R}^2_p
\]
be their Lagrangian product. A facet of \(K\) is either a \(q\)-facet,
coming from an edge of \(K_q\), or a \(p\)-facet, coming from an edge of
\(K_p\). We call a single facet a one-facet block, and an ordered adjacent pair
of facets of the same type a two-facet block.

\begin{theorem}[Finite enumeration for Lagrangian products]
\label{thm:lagrangian-product-finite-enumeration}
For a Lagrangian product \(K=K_q\times_L K_p\) of nondegenerate planar convex
polygons containing the origin in their interiors, \(c_{\mathrm{EHZ}}(K)\) is
obtained by maximizing \(Q_\sigma\) over the Haim--Kislev quadratic programs
associated to cyclic words of the form
\[
  \bigl([Q\mid QQ]\,[P\mid PP]\bigr)^k,
  \qquad k\in\{2,3\},
\]
where each \(Q\)-block is a one- or two-facet block in \(K_q\), each
\(P\)-block is a one- or two-facet block in \(K_p\), distinct blocks use
disjoint facets, and a two-facet block must use adjacent polygon edges. This is
an exhaustive candidate family for the capacity value; it may contain
overcounting or infeasible words, and it is not a classification of all
minimum-action simple Reeb orbits.
\end{theorem}

\begin{proof}
The Haim--Kislev quadratic program,
Theorem~\ref{thm:quadratic-program-algorithm-hk2019-hk-formula}, reduces
\(c_{\mathrm{EHZ}}(K)\) to a finite optimization over simple ordered facet
words. Thus it remains to find a finite subfamily that contains at least one
capacity-realizing word for a Lagrangian product.

On a \(q\)-facet of \(K_q\times_L K_p\), the Reeb vector is tangent to the
\(p\)-factor; on a \(p\)-facet, it is tangent to the \(q\)-factor. Hence a
simple orbit cannot contain three consecutive facets of the same type. If
three consecutive facets were \(q\)-facets, for example, the \(q\)-coordinate
would be constant through the two intermediate Reeb segments. The first
breakpoint would force this fixed point of \(K_q\) to lie in the intersection
of the first two \(K_q\)-edges, while the second breakpoint would force the
same point to lie in the intersection of the second and third \(K_q\)-edges.
For three distinct consecutive edges of a nondegenerate polygon these are the
two different endpoints of the middle edge, which is impossible for one fixed
point. The same argument applies to \(p\)-facets.

Therefore every same-type run in a cyclic facet word has length one or two,
and the maximal same-type runs alternate between \(q\)- and \(p\)-runs. A
two-facet run must use adjacent polygon edges, because consecutive facets in a
simple orbit share the breakpoint between the corresponding Reeb segments.
Since active Haim--Kislev words have distinct facet entries, distinct same-type
blocks use disjoint facets; otherwise the flattened cyclic word would repeat a
facet. A word with one \(q\)-block and one \(p\)-block cannot satisfy the
positive closure equation: the \(q\)-part would require a positive combination
of one normal, or of two adjacent edge normals, to vanish, and the same
obstruction holds in the \(p\)-part. Adjacent outward normals of a convex
polygon are not antiparallel. Thus \(k=1\) is excluded.

The EHZ capacity of a Lagrangian product is also the length of the shortest
closed \(K_p\)-billiard trajectory in \(K_q\), measured with the support
function of \(K_p\): Artstein--Avidan--Ostrover prove this for Lagrangian
products \cite[Theorem~2.13]{AAO2014}, and Rudolf records the nonsmooth
Minkowski-billiard form used here \cite[Theorem~1]{Rudolf2022}. The planar
billiard reduction of Bezdek and Bezdek
\cite[Theorem~1.1 and Lemma~2.4]{BezdekBezdek2009}, in the form used for
Minkowski billiards by Rudolf \cite[Theorem~1]{Rudolf2022}, implies that one
minimizing trajectory may be chosen with at most three bounce points. Under the
Reeb-orbit--billiard correspondence, a bounce point is exactly a maximal
interval of the Reeb orbit on which the \(q\)-coordinate is fixed and lies on
the active \(K_q\)-edge or adjacent edge pair; in the facet word this is one
\(q\)-block. The blocks alternate cyclically, so the number of \(p\)-blocks is
the same as the number of \(q\)-blocks. Hence \(k\leq 3\). Combining this with
\(k\neq 1\) gives \(k\in\{2,3\}\). Applying the Haim--Kislev quadratic program
to this alternating-block candidate family therefore gives the same minimum as
the full Lagrangian-product problem. This argument only proves that the
restricted family contains a capacity minimizer; it does not exclude additional
minimum-action simple Reeb orbits with more blocks.
\end{proof}

\begin{algorithm}[Lagrangian-product capacity enumeration]
\label{alg:lagrangian-product-capacity-enumeration}
For each factor, enumerate all one-facet blocks and all ordered adjacent
two-facet blocks. For \(k=2,3\), choose \(k\) non-overlapping \(q\)-blocks and
\(k\) non-overlapping \(p\)-blocks, order them cyclically in the alternating
pattern
\[
  \bigl([Q\mid QQ]\,[P\mid PP]\bigr)^k,
\]
flatten the blocks to a facet word, discard infeasible words, and solve the
corresponding Haim--Kislev quadratic program. The capacity is
\(1/(2Q_{\max})\), where \(Q_{\max}\) is the largest objective value among the
feasible words in this candidate family.
\end{algorithm}

\subsection{Use in later certificates}
\label{subsec:quadratic-program-algorithm-hk2019-implementation-status}

The later exact certificates use this section as a mathematical interface. The
HKO local-maximum certificate uses selected active words to build exact upper
branches. The rotated-pentagon certificate uses the specialized
Lagrangian-product enumeration from
Theorem~\ref{thm:lagrangian-product-finite-enumeration}, then checks the exact
algebraic inequalities needed to compare every surviving branch with the
displayed one.
